% Source: https://artofproblemsolving.com/community/c3h2128076_cool_problem_which_i_made
% Stable identity: aops-topic-2128076-q01
A parabola with a vertical axis and vertex $(-1,2)$ has an equation of the form
\[y=a(x+1)^2+2 .\]
Since $(1,3)$ is on the parabola, $3=a(1+1)^2+2=4a+2$, so $a=\dfrac{1}{4}$.

\begin{asy}
usepackage("amsmath");
size(5cm);
defaultpen(fontsize(10pt));
import graph;
real f(real x){return 0.25*(x+1)^2+2;}
draw((-4.6,0)--(2.6,0),Arrow(6)); draw((0,-0.4)--(0,5.2),Arrow(6));
label("$x$",(2.6,0),E); label("$y$",(0,5.2),N);
draw(graph(f,-4.3,2.3),blue+linewidth(1));
draw((-1,0)--(-1,4.9),gray+dashed);
dot((-1,2)); dot((1,3)); dot((-3,3),gray);
label("$V(-1,2)$",(-1,2),SW,fontsize(9pt));

label("$(-3,3)$",(-3,3),SW,gray+fontsize(9pt));
label("$x=-1$",(-1,4.9),N,gray+fontsize(8pt));
label("$(1,3)$",(1,3),SE,fontsize(9pt));
label("\mbox{}",(min(currentpicture,user=true)+max(currentpicture,user=true))/2,black); // keeps the text after the figure black when LaTeX typesets the labels inline
\end{asy}

Expanding,
\[y=\dfrac{1}{4}\left(x^2+2x+1\right)+2=\boxed{\dfrac{1}{4}x^2+\dfrac{1}{2}x+\dfrac{9}{4}}.\]
Check: here $a=\dfrac{1}{4}$ and $b=\dfrac{1}{2}$, so $2a=\dfrac{1}{2}=b$ and the vertex is at $x=-\dfrac{b}{2a}=-1$; also $y(-1)=\dfrac{1}{4}-\dfrac{1}{2}+\dfrac{9}{4}=2$; also $y(1)=\dfrac{1}{4}+\dfrac{1}{2}+\dfrac{9}{4}=3$. By symmetry the parabola also passes through $(-3,3)$.
